\section{Introduction}

Cycles are a powerful concept in computer science and mathematics,
representing unbounded lists that repeat a finite sequence of elements. When
we integrate cycles, we obtain a list of cumulative sums with some unique
properties that we will explore in this article.

\begin{equation*}
\begin{aligned}
L &= [l_0, l_1, l_2, \ldots, l_{n-1}]  \mid &l_n &\in \mathbb{S}, L \in \mathbb{L}
\end{aligned}
\end{equation*}

\begin{equation*}
\begin{aligned}
\operatorname{Cycle}(L)               &= [l_0, l_1, l_2, \ldots, l_{n-1}, l_0, l_1, \ldots] \\
                              &= [v_0, v_1, v_2, \dots] \mid &v_i &= L[i \operatorname{mod} n] \\
\operatorname{Integral}(L, init)      &= [y_0, y_1, y_2, \ldots, y_{n-1}] \mid &y_k &= \sum_{i=0}^{k} x_i + init \\
\operatorname{CycleIntegral}(L, init) &= [w_0, w_1, w_2, \ldots] \mid          &w_k &= \sum_{i=0}^{k} \operatorname{Cycle}(L)_i + init
\end{aligned}
\end{equation*}

In this article, we present a discrete definition of Cycle Integral over
finite integer lists, defined recursively and verified some of its properties
using the Stainless system. Our approach follows a zero-prior-knowledge
philosophy, building on a previously verified foundation for recursive list
and integral structures and summation. The result is a verified, from-scratch
implementation of the cycle integral suitable as a foundation for
higher-level numeric reasoning over unbounded lists.

This article verifies:

\begin{itemize}
  \item Two equivalent definitions: recursive and modulo-based ---
    Subsections~3.1--3.3
  \item Core properties: next position, same difference after cycle, sum of
    mod values, strictly increasing, positivity --- Subsections~4.1--4.5
  \item Persistent and periodic properties: how a fixed cycle integral's
    residues behave forever, and how it advances across full periods ---
    Subsections~5.1--5.6
  \item Deriving new cycle integrals: expansion, index shifts, rotation,
    survivor filtering, and merge-based reconstruction ---
    Subsections~6.1--6.10
\end{itemize}

\subsection{Related work}

Lean's Mathlib provides a general formal theory of periodic functions. It
proves that a periodic function remains periodic under integer multiples of a
period, and that a finite sum of periodic functions is periodic~\cite{leanmathlibperiodic}.
Mathlib also represents a finite cycle as a list modulo cyclic
rotation~\cite{leanmathlibcycles}. These results give formal context for the
finite period and shift structure used by the present construction.

This article studies a different object: the cumulative integral of a
concrete periodic integer list. Its central verified results are the
equivalence of a recursive integral and a quotient--remainder closed form,
together with the resulting step, full-period, residue, and reconstruction
properties. The existing periodicity and finite-cycle developments therefore
enrich the setting without standing in for this two-presentation
cycle-integral proof.
